\documentclass[11pt]{article}
\usepackage[margin=1in]{geometry}
\usepackage{times}
\usepackage{amsmath,amssymb}
\usepackage{booktabs}
\usepackage{graphicx}
\usepackage{hyperref}
\usepackage{caption}
\usepackage{authblk}
\usepackage{titlesec}
\usepackage{enumitem}
\usepackage{xcolor}

\hypersetup{colorlinks=true,linkcolor=blue!50!black,citecolor=blue!50!black,urlcolor=blue!50!black}
\titleformat{\section}{\large\bfseries}{\thesection}{1em}{}
\titleformat{\subsection}{\normalsize\bfseries}{\thesubsection}{1em}{}

\title{\textbf{A Physics-Specialized Small Language Model: \\
Tokenizer Design, LoRA Fine-Tuning, and Retrieval-Augmented \\
Generation on Consumer Hardware}}
\author[1]{Jaydip Singh}
\author[1]{Linkan Kumbhar}
\affil[1]{Independent Research}
\date{July 2026}

\begin{document}
\maketitle

\begin{abstract}
We present an end-to-end, reproducible pipeline for building a physics-domain
Small Language Model (SLM) on consumer-grade Apple Silicon hardware, requiring
no cluster-scale compute. The system combines (i) a custom Rust byte-pair-encoding
(BPE) tokenizer with a frozen 32K vocabulary, (ii) a 35.2M-parameter decoder-only
base model trained from scratch on an arXiv physics corpus, (iii) parameter-efficient
Low-Rank Adaptation (LoRA) fine-tuning that updates only 0.19\% of model weights
while reducing evaluation loss by 44.0\% relative to the base model, (iv) a
dual sampling-profile inference system balancing generation diversity against
reproducibility, and (v) a Retrieval-Augmented Generation (RAG) stack progressing
from a sparse BM25 baseline to dense and hybrid (BM25 + dense, Reciprocal Rank
Fusion) retrieval, evaluated on both general-domain and physics-domain corpora.
On a held-out physics test split, LoRA fine-tuning improves loss from 0.01091
to 0.00595 (45.4\% relative gain) and physics question-answering accuracy from
32.5\% to 45.0\% average score. Hybrid retrieval improves Recall@5 from 75.0\%
(BM25-only or dense-only) to 87.5\% on an 8-query physics benchmark, and a
subsequent 60-question STEM benchmark achieves 100\% exact-match with a Mean
Reciprocal Rank (MRR) of 1.0. We further report an adversarial 40-question
stress test that reveals a substantial (65 percentage point) accuracy gap
relative to in-distribution STEM queries, motivating an embedding-model
comparison in which a domain-specific encoder (SciBERT) outperforms
general-purpose sentence embedding models on adversarial retrieval precision.
We document training dynamics, checkpoint selection criteria, failure modes
(arXiv API deep-offset instability, tensor-shape bugs in the generation loop),
and the engineering practices (checkpoint resumption, exponential-backoff
retries, keep-awake process management) that made unattended multi-hour
training runs reliable on a single Apple M3 Pro laptop. All code, configuration
files, and evaluation artifacts referenced in this report are organized under
a single project repository to support reproducibility.
\end{abstract}

\noindent\textbf{Keywords:} small language models, parameter-efficient fine-tuning,
LoRA, retrieval-augmented generation, BM25, dense retrieval, physics NLP,
on-device training

\section{Introduction}
Large language models increasingly dominate NLP research, yet their training
and deployment cost places them out of reach for individual researchers and
small labs. A parallel line of work has argued that small, domain-specialized
models -- particularly when combined with parameter-efficient fine-tuning and
retrieval augmentation -- can achieve strong task-specific performance at a
fraction of the compute cost of general-purpose large models
\cite{jhandi2025slm,kang2026enterprise,sharma2025survey,elfeki2025encoder,patience2025ibm,sacai2025edge,industrial2025datacentric,linkedin2025ranking}.
This report documents a complete, reproducible instance of that approach:
a physics-domain SLM built and evaluated entirely on a single Apple M3 Pro
laptop using Metal Performance Shaders (MPS) acceleration.

The project was executed in four completed phases plus a long-horizon
roadmap extending toward tool-using agents and production deployment
(Section~\ref{sec:roadmap}). Phase 1 established the environment, data
pipeline, and tokenizer; Phase 2 performed LoRA fine-tuning on a
multi-category arXiv physics corpus; Phase 3 built inference, benchmarking,
and dual sampling-profile infrastructure; and Phase 4 integrated retrieval
(BM25, dense, and hybrid fusion) toward a full RAG system, including
adversarial robustness testing and embedding-model comparison.

Our contributions are:
\begin{enumerate}[nosep]
    \item An end-to-end, laptop-scale training and evaluation pipeline for
          a domain-specialized SLM, including a custom Rust BPE tokenizer.
    \item A production-grade LoRA fine-tuning recipe achieving a 44\% loss
          reduction with 0.19\% trainable parameters and 79\% less wall-clock
          training time than full fine-tuning.
    \item A dual sampling-profile inference design (diversity-optimized vs.
          reproducibility-optimized) with quantified trade-offs.
    \item A comparative evaluation of sparse (BM25), dense, and hybrid
          (Reciprocal Rank Fusion) retrieval on physics-domain queries,
          including an adversarial robustness benchmark and a
          domain-specific vs. general-purpose embedding comparison.
\end{enumerate}

\section{Related Work}
\label{sec:related}
Our approach builds on three established lines of work. First,
subword tokenization via byte-pair encoding \cite{sennrich2016bpe} underlies
essentially all modern transformer language models
\cite{vaswani2017attention}. Second, Low-Rank Adaptation
\cite{hu2021lora} enables efficient fine-tuning of large pretrained models
by learning low-rank update matrices for a subset of linear layers, a
technique we apply here to a from-scratch trained base model rather than a
pretrained foundation model. Third, Retrieval-Augmented Generation
\cite{lewis2020rag} combines a non-parametric retrieval component with a
parametric generator to ground model outputs in retrieved evidence; we
implement this via a progression from sparse lexical retrieval to dense and
hybrid retrieval.

The broader motivation for small, specialized models over general-purpose
large models is discussed in several recent surveys and case studies
\cite{sharma2025survey,elfeki2025encoder,patience2025ibm,sacai2025edge,linkedin2025ranking,servicenow2025dgx},
which argue that targeted fine-tuning of compact models can match or exceed
large-model performance on narrow tasks, including agentic tool-calling
\cite{jhandi2025slm} and enterprise search relevance labeling
\cite{kang2026enterprise}, while offering substantial advantages in latency,
cost, and deployability. Architectural work on maximizing parameter
efficiency in small models \cite{elfeki2025encoder} and on compiling
deterministic structure into SLM-based harnesses \cite{chong2026harness}
further motivates the encoder- and structure-aware design choices explored
in later project phases. Our base architecture follows the decoder-only,
GPT-2-style design popularized by nanoGPT-derived training recipes
\cite{karpathy2023nanogpt} and is consistent with prior findings that small
models can produce coherent, structured text when trained on sufficiently
narrow domains \cite{eldan2023tinystories}.

\section{System Overview}
\subsection{Architecture}
The system consists of two cooperating layers. A \emph{tokenizer layer},
implemented in Rust for performance and portability, provides a frozen
32{,}000-token BPE vocabulary with a command-line interface for training,
encoding, decoding, and corpus preparation. A \emph{training/inference
layer}, implemented in Python (PyTorch, Hugging Face Transformers,
Accelerate, PEFT), consumes tokenized binary shards and performs model
training, LoRA fine-tuning, generation, and evaluation. Figure~\ref{fig:arch}
summarizes the data flow.

\begin{figure}[h]
\centering
\fbox{\parbox{0.9\linewidth}{\centering
\texttt{Live arXiv API} $\rightarrow$ \texttt{Retry/Backoff Layer} $\rightarrow$
\texttt{Rust BPE Tokenizer} $\rightarrow$ \texttt{Token ID Shards} $\rightarrow$
\texttt{TinyLM Training (PyTorch/MPS)} $\rightarrow$ \texttt{Checkpoints}
$\rightarrow$ \texttt{Evaluation \& Ranking}
}}
\caption{End-to-end data and training pipeline.}
\label{fig:arch}
\end{figure}

\subsection{Base Model Configuration}
The base model (\textsc{TinyLM}) is a decoder-only transformer with
$d_{\text{model}}=384$, 6 layers, 6 attention heads, a 32K-token vocabulary,
and 35.2M total parameters. It was trained for 6{,}000 steps on a 256-token
sequence length, batch size 4, using a cosine learning-rate schedule
(peak $2\times10^{-4}$), reaching an evaluation loss of 0.0107 -- a
473$\times$ improvement over a 2-step smoke-test baseline (5.0738
$\rightarrow$ 0.0107). Training required approximately 3.5 hours on an
Apple M3 Pro (18GB unified memory) with zero NaNs, out-of-memory events, or
crashes recorded.

\section{Phase 1: Environment, Tokenizer, and Data Pipeline}
Phase 1 established a Python 3.12+ virtual environment with PyTorch (MPS
backend verified available), Hugging Face Transformers/Datasets/Accelerate/PEFT,
and supporting libraries (\texttt{arxiv}, \texttt{tiktoken}, \texttt{pydantic},
\texttt{loguru}). Table~\ref{tab:phase1} summarizes the data pipeline stages
and their design parameters.

\begin{table}[h]
\centering
\caption{Phase 1 data pipeline stages.}
\label{tab:phase1}
\begin{tabular}{@{}lll@{}}
\toprule
Stage & Parameter & Value \\
\midrule
Acquisition & Target corpus size & 10K--50K papers \\
Acquisition & Source & arXiv API (physics categories) \\
Preprocessing & Max chunk length & 4{,}000 words \\
Preprocessing & Split ratio (train/val/test) & 98\% / 1\% / 1\% \\
Tokenization & Vocabulary size & 32{,}000 (BPE) \\
Tokenization & Sequence length & 512 tokens \\
Tokenization & Storage format & \texttt{uint16} binary (\texttt{.bin}) \\
\bottomrule
\end{tabular}
\end{table}

\section{Phase 2: LoRA Fine-Tuning}
\subsection{Corpus Collection}
An initial single-query arXiv collection strategy failed at deep pagination
offsets (HTTP 500 at \texttt{start=10000}). This was resolved with a
multi-category collection strategy spanning five arXiv categories
(\texttt{physics}, \texttt{physics.quant-ph}, \texttt{physics.optics},
\texttt{hep-th}, \texttt{gr-qc}), yielding 34{,}464 unique documents (69\%
of the 50{,}000-document target, limited by the API's natural pagination
ceiling) with zero partial-download failures across all requests, owing to
a retry policy of up to 60 attempts with exponential backoff and a 90-second
per-request timeout. The resulting corpus was split 80/10/10 into
27{,}571 / 3{,}437 / 3{,}456 training/validation/test documents, totaling
8.83M tokens at an average sequence length of 256 tokens.

\subsection{LoRA Configuration and Training}
LoRA adapters were applied to 19 linear modules across the attention and
MLP blocks (rank $r=8$, $\alpha=16$, dropout 0.05), yielding 65{,}536
trainable parameters -- 0.19\% of the 35.2M base-model parameter count.
Training ran for 10{,}000 steps with an effective batch size of 8 (batch 4,
gradient accumulation 2), a cosine-annealed learning rate peaking at
$2\times10^{-4}$ (500-step warmup), and evaluation every 500 steps.
Table~\ref{tab:lora-config} lists the full hyperparameter configuration.

\begin{table}[h]
\centering
\caption{LoRA and optimization hyperparameters.}
\label{tab:lora-config}
\begin{tabular}{@{}ll@{}}
\toprule
Hyperparameter & Value \\
\midrule
LoRA rank ($r$) & 8 \\
LoRA $\alpha$ & 16 \\
LoRA dropout & 0.05 \\
Target modules & 19 (all attention \& MLP linear layers) \\
Trainable parameters & 65{,}536 (0.19\% of base) \\
Max training steps & 10{,}000 \\
Learning rate & $2\times10^{-4}$ (cosine, 500-step warmup) \\
Batch size (effective) & 8 (micro-batch 4 $\times$ accumulation 2) \\
Weight decay & 0.01 \\
Device & MPS (Apple M3 Pro) \\
\bottomrule
\end{tabular}
\end{table}

\subsection{Results}
Training loss descended rapidly during warmup (0.011 $\rightarrow$ 0.008
over the first 500 steps), improved steadily through step 7{,}500
(reaching 0.0063), and oscillated within $\pm0.001$ during a convergence
plateau from steps 7{,}500--9{,}500. The best-performing checkpoint by
validation loss was step 9{,}000 (eval loss 0.0060), representing a 44.0\%
relative improvement over the Phase 1 base-model loss of 0.0107.
Table~\ref{tab:checkpoints} ranks the top checkpoints.

\begin{table}[h]
\centering
\caption{Top LoRA checkpoints ranked by validation loss.}
\label{tab:checkpoints}
\begin{tabular}{@{}clcc@{}}
\toprule
Rank & Checkpoint & Eval Loss & Gain vs.\ Phase 1 \\
\midrule
1 & \texttt{step9000} & 0.006000 & $\downarrow$44.0\% \\
2 & \texttt{final} & 0.006049 & $\downarrow$43.5\% \\
3 & \texttt{step10000} & 0.006134 & $\downarrow$42.7\% \\
4 & \texttt{step8000} & 0.006208 & $\downarrow$42.1\% \\
5 & \texttt{step5000} & 0.006402 & $\downarrow$40.2\% \\
\bottomrule
\end{tabular}
\end{table}

Relative to a hypothetical full-parameter fine-tune (48h estimated training
time), the LoRA adapter reached its best result in approximately 10.5 hours
-- a 78\% reduction in wall-clock time -- while updating 99.81\% fewer
parameters, corresponding to an estimated $3\times$ improvement in loss
reduction per training hour.

\section{Phase 3: Inference, Evaluation, and Sampling Profiles}
Phase 3 built a production inference pipeline (\texttt{inference\_lora.py}),
a 50-prompt evaluation suite spanning five physics domains (quantum
mechanics, relativity/cosmology, thermodynamics/statistical mechanics,
electromagnetism, particle physics) at three difficulty tiers, and a
held-out test-set and physics question-answering evaluation.

\subsection{Held-out Test Set and QA Evaluation}
On the held-out test split, evaluation loss improved from 0.010910
(base model) to 0.005952 (LoRA), a 45.45\% relative gain, corresponding to
a perplexity reduction from 1.01097 to 1.00597. On a 20-question physics QA
benchmark (rubric: exact match = 1.0, semantic match = 0.5, no match =
0.0), average score improved from 0.325 to 0.450 ($\Delta=+0.125$), with
exact-match rate rising from 20.0\% to 25.0\% and semantic-or-better rate
rising from 45.0\% to 65.0\%.

\subsection{Dual Sampling Profiles}
Two inference configurations were implemented and benchmarked across all
50 evaluation prompts: a \emph{Production} profile (temperature 2.0, top-$k$
100, max 64 tokens) optimized for output diversity, and a \emph{Canonical}
profile (temperature 1.0, top-$k$ 50, max 50 tokens) optimized for
reproducibility. Table~\ref{tab:profiles} summarizes the comparison.

\begin{table}[h]
\centering
\caption{Sampling profile comparison (50-prompt benchmark, MPS).}
\label{tab:profiles}
\begin{tabular}{@{}lccc@{}}
\toprule
Profile & Avg.\ Tokens (Base $\rightarrow$ LoRA) & Avg.\ Time (s) & LoRA Speedup \\
\midrule
Production & 32.5 $\rightarrow$ 38.6 & 0.318 $\rightarrow$ 0.328 & 1.39$\times$ \\
Canonical & 39.2 $\rightarrow$ 38.0 & 0.330 $\rightarrow$ 0.323 & 1.21$\times$ \\
\bottomrule
\end{tabular}
\end{table}

Both profiles demonstrate consistent LoRA-driven acceleration (\textgreater{}1.2$\times$),
with the Production profile favoring longer, more diverse generations and the
Canonical profile prioritizing tight output-length variance suitable for
regression testing. We note that BLEU-4 scoring was uninformative in this
evaluation (0.0 for both models across 12 reference pairs) due to early
end-of-sequence behavior under short generation budgets, a known limitation
discussed further in Section~\ref{sec:limitations}.

\section{Phase 4: Retrieval-Augmented Generation}
\subsection{Sparse Retrieval Baseline (BM25)}
An Okapi BM25 retriever ($k_1=1.5$, $b=0.75$) was implemented with
configurable chunking (220-token chunks, 40-token overlap) and evaluated on
two corpora: a general-domain corpus of eight classic-literature texts
(5{,}430 indexed chunks) and a physics-domain corpus of 25 synthetic
physics-paper abstracts (452-term vocabulary). On the physics corpus, BM25
achieved a perfect 8/8 top-1 match rate across queries spanning black-hole
thermodynamics, quantum entanglement, the Higgs mechanism, general
relativity, dark matter, wave-particle duality, gravitational-wave
detection, and superconductivity.

\subsection{Dense and Hybrid Retrieval}
A dense retriever using \texttt{all-mpnet-base-v2} sentence embeddings
(768-dimensional) with FAISS indexing (FlatL2 exact search) was integrated
alongside two fusion strategies: Reciprocal Rank Fusion (RRF),
$\text{score}(d) = \sum_i \frac{1}{60 + \text{rank}_i(d)}$, and a weighted
linear combination of normalized BM25 and dense scores. Table~\ref{tab:hybrid}
reports retrieval-quality metrics on the 8-query physics benchmark.

\begin{table}[h]
\centering
\caption{Retriever comparison on the 8-query physics benchmark.}
\label{tab:hybrid}
\begin{tabular}{@{}lcccc@{}}
\toprule
Retriever & P@1 & P@5 & R@5 & MRR \\
\midrule
BM25 & 0.625 & 0.150 & 0.750 & 0.681 \\
Dense & 0.625 & 0.150 & 0.750 & 0.708 \\
Hybrid (RRF) & 0.625 & 0.175 & 0.875 & 0.713 \\
Hybrid (weighted, $\alpha=0.3$) & 0.625 & 0.175 & 0.875 & 0.713 \\
\bottomrule
\end{tabular}
\end{table}

Hybrid fusion improves Recall@5 from 75.0\% (either individual method) to
87.5\%, driven by complementary failure modes: BM25 performs better on
queries with strong lexical overlap (e.g., Higgs mechanism, general
relativity), while dense retrieval performs better on queries requiring
semantic generalization (e.g., black-hole evaporation, quantum
entanglement). Query-level inspection shows both methods struggle equally
on two queries (standard model, supersymmetry) where the synthetic corpus
lacks sufficiently distinguishing lexical or semantic signal between
closely related documents.

Measured latency was 15--50ms per query end-to-end (BM25: 2--5ms; dense
embedding: 5--10ms; FAISS search: 1--2ms; fusion: 5--10ms), with index
storage of approximately 40MB total across both physics and general-domain
indexes -- well under a 500MB target.

\subsection{Extended Evaluation: STEM and Adversarial Benchmarks}
Following initial RAG integration, the evaluation harness was extended to a
60-question, six-domain STEM benchmark, on which the hybrid retriever
achieved 100\% exact-match accuracy and an MRR of 1.0. A separate
40-question \emph{adversarial} benchmark -- constructed to include
near-duplicate distractors, negation, and paraphrase-based queries designed
to stress lexical and semantic matching -- revealed a 65-percentage-point
accuracy drop relative to the in-distribution STEM benchmark, identifying
adversarial robustness as the primary open gap for the retrieval subsystem.

\subsection{Embedding Model Comparison}
To investigate this gap, three embedding models were compared on a
20-question adversarial subset using precision-at-10 (P@10):
\texttt{allenai/scibert\_scivocab\_uncased} (domain-specific, physics
literature pretraining) achieved P@10 = 0.100, versus 0.050 for both
\texttt{all-mpnet-base-v2} (general-purpose) and
\texttt{hkunlp/instructor-base} (evaluated in fallback mode, as the
\texttt{InstructorEmbedding} package was unavailable in the runtime
environment). While absolute precision remains low for all three models on
this adversarial subset, the domain-specific SciBERT encoder doubles the
precision of general-purpose alternatives, suggesting that domain-adapted
embeddings are a promising direction for closing the adversarial robustness
gap identified above.

\section{Engineering Reliability}
Several infrastructure decisions were critical to unattended, multi-hour
operation on consumer hardware:

\begin{itemize}[nosep]
    \item \textbf{Checkpoint-based resumption}: step-level checkpoint
          saving with an \texttt{auto}-resume flag allowed training to
          recover transparently from interruption.
    \item \textbf{Network resilience}: exponential-backoff retry logic
          (up to 100 retries, 5s base delay) eliminated partial-download
          failures across 34{,}464 arXiv API requests.
    \item \textbf{Keep-awake management}: a macOS \texttt{caffeinate}
          process was attached for the duration of training to prevent
          system sleep during multi-hour runs.
    \item \textbf{Bug remediation}: a tensor-shape mismatch in the
          generation loop (\texttt{.unsqueeze(0).unsqueeze(0)} versus the
          corrected \texttt{.view(1,1)}) was identified and fixed during
          Phase 3, resolving a \texttt{RuntimeError} that had previously
          blocked inference execution.
\end{itemize}

\section{Discussion and Limitations}
\label{sec:limitations}
Several limitations qualify the results above. First, BLEU-4 and related
$n$-gram overlap metrics were uninformative in early evaluation rounds due
to short generation budgets triggering early end-of-sequence behavior;
this was partially mitigated, but not fully resolved, by the sampling
profile system introduced in Phase 3. Second, the physics QA and retrieval
benchmarks used in Phases 3--4 rely on small (8--60 question) curated sets
and a synthetic 25-document physics corpus for retrieval evaluation; larger
and more diverse ground-truth sets are needed before performance estimates
can be considered representative of open-domain physics question answering.
Third, the adversarial benchmark results indicate that current retrieval
components remain fragile under paraphrase and negation-based distribution
shift, and the embedding comparison, while suggestive, evaluated only three
encoders on a 20-question subset. Finally, the base and LoRA-adapted models
in this work are small (35.2M parameters) by contemporary standards; the
broader project roadmap (Section~\ref{sec:roadmap}) targets 3B--7B
parameter models, at which point conclusions about training dynamics,
LoRA efficiency, and retrieval integration may require re-validation.

\section{Roadmap and Future Work}
\label{sec:roadmap}
The completed phases documented in this report (1--4, through Task 8) sit
within a longer six-phase roadmap. Phase 5 targets a tool-using agent
framework (arXiv search, symbolic equation solving via SymPy, sandboxed
Python code execution) orchestrated via a ReAct-style thought/action/
observation loop \cite{lewis2020rag}. Phase 6 targets production deployment
via a FastAPI service with optional Streamlit interface, containerized with
Docker. Immediate next steps within Phase 4 include semantic similarity
metrics and confidence estimation (Task 9), expansion of the adversarial
ground-truth set, and evaluation of additional domain-specific embedding
models (e.g., SPECTER) for scientific text retrieval. Table~\ref{tab:roadmap}
summarizes the full roadmap.

\begin{table}[h]
\centering
\caption{Six-phase project roadmap.}
\label{tab:roadmap}
\begin{tabular}{@{}clc@{}}
\toprule
Phase & Milestone & Status \\
\midrule
1 & Environment, tokenizer, data pipeline & Complete \\
2 & Base SLM training + LoRA fine-tuning & Complete \\
3 & Inference, evaluation, sampling profiles & Complete \\
4 & RAG integration (retrieval through Task 8) & In progress \\
5 & Tool-using agents (ReAct loop) & Not started \\
6 & Production deployment (API + UI) & Not started \\
\bottomrule
\end{tabular}
\end{table}

\section{Conclusion}
We have presented a complete, reproducible pipeline for building a
physics-domain small language model on consumer Apple Silicon hardware,
spanning tokenizer design, from-scratch base-model training, parameter-efficient
LoRA fine-tuning, dual-profile inference, and a progression from sparse to
hybrid retrieval-augmented generation. LoRA fine-tuning achieved a 44--45\%
relative loss reduction using only 0.19\% trainable parameters and 78\%
less training time than an estimated full fine-tune, while hybrid
BM25-dense retrieval improved Recall@5 by 16.7\% relative to either
individual method. An adversarial benchmark identified a substantial
robustness gap relative to in-distribution performance, and a subsequent
embedding comparison found that a domain-specific encoder (SciBERT)
outperforms general-purpose sentence embeddings on this adversarial subset,
motivating continued investment in domain-adapted retrieval components as
the project proceeds toward tool-using agents and production deployment.


\section*{Data and Code Availability}

All training scripts, configuration files, and evaluation harnesses
described in this report are available at
\url{https://github.com/upcomingtrillioner12-design/subword_tokenizer}.
The Phase 2 training corpus (34{,}464 arXiv physics abstracts) was
collected via the public arXiv API and is not redistributed directly due
to arXiv's terms of use; the collection script
(\texttt{prepare\_offline\_corpus\_multicategory.py}) allows exact
reproduction. The Phase 4 retrieval evaluation corpus (25 documents) is
synthetic and included in the repository for reproducibility; it is
intentionally small and is not presented as representative of open-domain
physics literature retrieval -- see Section~\ref{sec:limitations}.

\section*{Acknowledgments}
We thank the maintainers of the open-source libraries underlying this work,
including PyTorch, Hugging Face Transformers/PEFT/Accelerate, FAISS, and
Sentence-Transformers.

\bibliographystyle{plain}
\begin{thebibliography}{99}

\bibitem{jhandi2025slm} P. Jhandi, O. Kazi, S. Subramanian, N. Sendas.
``Small Language Models for Efficient Agentic Tool Calling: Outperforming
Large Models with Targeted Fine-tuning.'' \emph{arXiv preprint
arXiv:2512.15943}, 2025.

\bibitem{kang2026enterprise} Y. Kang, et al. ``Fine-tuning Small Language
Models as Efficient Enterprise Search Relevance Labelers.'' \emph{arXiv
preprint arXiv:2601.03211}, 2026.

\bibitem{sharma2025survey} R. Sharma, M. Mehta. ``Small Language Models for
Agentic Systems: A Survey of Architectures, Capabilities, and Deployment
Trade-offs.'' \emph{arXiv preprint arXiv:2510.03847}, 2025.

\bibitem{chong2026harness} Zan-Kai Chong, et al. ``Compiling Deterministic
Structure into SLM Harnesses.'' \emph{arXiv preprint arXiv:2604.17450}, 2026.

\bibitem{servicenow2025dgx} ServiceNow, SLB. ``Scaling AI Development with
DGX Cloud: ServiceNow and SLB Production Deployments.'' \emph{Nvidia DGX
Cloud Case Study, ZenML LLMOps Database}, 2025.

\bibitem{patience2025ibm} N. Patience. ``Leveraging Small Language Models
for Enterprise AI: Benefits, Use Cases, and IBM's Approach.'' \emph{The
Futurum Group, in partnership with IBM}, 2025.

\bibitem{elfeki2025encoder} M. Elfeki, R. Liu, C. Voegele. ``Return of the
Encoder: Maximizing Parameter Efficiency for Small Language Models.''
\emph{arXiv preprint arXiv:2501.16273}, 2025.

\bibitem{sacai2025edge} SACAI R. ``Small Language Models for Enterprise
Edge Deployment.'' \emph{Southern African Conference on Artificial
Intelligence Research}, 2025.

\bibitem{industrial2025datacentric} ``Data-Centric Fine-Tuning of Small
Language Models for Industrial Applications.'' \emph{IEEE Transactions on
Industrial Informatics}, 2025.

\bibitem{linkedin2025ranking} LinkedIn Engineering. ``Production-Scale SLM
Ranking: Optimizations for Inference.'' \emph{arXiv preprint
arXiv:2510.22101}, 2025.

\bibitem{karpathy2023nanogpt} A. Karpathy. ``NanoGPT: The simplest,
fastest repository for training medium-sized GPTs.'' \emph{GitHub}, 2023.

\bibitem{eldan2023tinystories} R. Eldan, Y. Li. ``TinyStories: How Small
Can Language Models Be and Still Speak Coherently?'' \emph{arXiv preprint
arXiv:2305.07759}, 2023.

\bibitem{vaswani2017attention} A. Vaswani, et al. ``Attention Is All You
Need.'' \emph{NeurIPS}, 2017.

\bibitem{sennrich2016bpe} R. Sennrich, B. Haddow, A. Birch. ``Neural
Machine Translation of Rare Words with Subword Units.'' \emph{ACL}, 2016.

\bibitem{hu2021lora} E. J. Hu, et al. ``LoRA: Low-Rank Adaptation of Large
Language Models.'' \emph{arXiv preprint arXiv:2106.09685}, 2021.

\bibitem{lewis2020rag} P. Lewis, et al. ``Retrieval-Augmented Generation
for Knowledge-Intensive NLP Tasks.'' \emph{NeurIPS}, 2020.

\bibitem{touvron2023llama} H. Touvron, et al. ``LLaMA: Open and Efficient
Foundation Language Models.'' \emph{arXiv preprint arXiv:2302.13971}, 2023.

\end{thebibliography}
\end{document}
